\documentclass[11pt]{article}
\usepackage[margin=1in]{geometry}
\usepackage{amsmath,amssymb,amsthm,mathtools,bm}
\usepackage{microtype}
\usepackage[T1]{fontenc}
\usepackage{lmodern}
\usepackage{booktabs,enumitem}
\usepackage[hidelinks]{hyperref}
\setlist{nosep}
\newtheorem{theorem}{Theorem}[section]
\newtheorem{corollary}[theorem]{Corollary}
\newtheorem{proposition}[theorem]{Proposition}
\theoremstyle{definition}
\newtheorem{definition}[theorem]{Definition}
\theoremstyle{remark}
\newtheorem{remark}[theorem]{Remark}

\newcommand{\AU}{\mathrm{AU}}
\newcommand{\Msun}{M_\odot}

\title{\textbf{Synopsis: Solar-System Null Test of the Hysterical Response}\\[0.35em]
\large What the bound says (without the orbital algebra)}
\author{Andrew Bruce Boyd}
\date{September 2026}

\begin{document}
\maketitle
\noindent\textit{Computational note.}
Proofs, exploratory experiments, and paper writing were assisted by generative AI.
Mathematical theorems stated in this paper are formally verified in Lean~4 in the
companion specifications, as faithful ports of the TeX arguments at the abstractions
recorded there, except results explicitly tagged as \emph{literature hypotheses}
(standard theorems cited from the literature and not re-proved here) or
\emph{physical inputs} (modeling assumptions outside mathematics).

\begin{abstract}
This is a short guide to the Solar-System null test of the Hysterical Response.
It assumes the response-theory and galactic-modelling papers.
Each section says \emph{why} the claim matters, then states the result.
Orbital algebra and observational detail live in the full paper.
\end{abstract}

\section{Why a Solar-System null test}
Galactic modelling supplies $g_{\mathrm{resp}}=CQ/r$ for response-supported systems.
That does not force the Sun into the same regime: the Solar amplitude $V_\odot^{2}=CQ_\odot$ is a separate, system-dependent quantity.
The question here is only how large $V_\odot$ can be before Saturn-scale ephemerides object.

\section{Point-source specialization}
Outside a localized Solar source, the same far-field law reads
\[
  g_{\mathrm{resp}}=\frac{V_\odot^{2}}{r},
  \qquad
  \epsilon(r)=\frac{V_\odot^{2}r}{GM_\odot}.
\]
The fractional correction grows with radius, so outer planets bite hardest.
An equivalent enclosed mass is $M_{\mathrm{resp}}(<r)=V_\odot^{2}r/G$.

\section{Perihelion signature}
A purely radial $1/r$ acceleration produces a retrograde perihelion drift.

\begin{proposition}[Perihelion shift of a logarithmic response]
\[
\Delta\varpi
=-\frac{2\pi V_\odot^{2} a}{GM_\odot}
\frac{\sqrt{1-e^{2}}}{1+\sqrt{1-e^{2}}}
\]
per orbit, reducing to $-\pi V_\odot^{2}a/GM_\odot$ for $e\to0$.
\end{proposition}

(The Gauss reduction to the true-anomaly integrand is a literature input; the remaining integral algebra is Lean-checked.)

\section{Saturn ceiling}
Translating the Pitjeva--Pitjev Saturn perihelion uncertainty $4.7\times10^{-6}$ arcsec\,yr$^{-1}$ through that formula yields
\[
  V_\odot<0.141\,\mathrm{m\,s^{-1}}
\]
at the quoted one-sigma level.
An enclosed-mass consistency scale from the same literature sits near $V_\odot\lesssim0.086\,\mathrm{m\,s^{-1}}$.

\section{Benchmark inside the window}
For $V_\odot=0.01\,\mathrm{m\,s^{-1}}$, Saturn sees
$g_{\mathrm{resp}}\sim7\times10^{-17}\,\mathrm{m\,s^{-2}}$ and
$\dot\varpi_{\mathrm{resp}}\sim-2.4\times10^{-6}$ arcsec per century---below the quoted limit.
Cassini $\gamma$ is an independent relativistic consistency condition, not a direct $V_\odot$ measurement.
A Saturn-centric amplitude $V_{\mathrm{Sat}}$ needs a dedicated Grand Finale re-fit.

\section{What is and is not claimed}
The far-field shape is inherited.
The perihelion formula is a theorem of the stated Gauss integrand.
The numerical ceiling is an observational translation, not a detection.
A nonzero subcritical window remains open; closing it requires explicit $V_\odot^{2}/r$ ephemeris fits.

\end{document}
